\begin{definition}[Three plaquette coordinates from directed bond labels]%
    \label{def:peps_regular_three_plaquette_coordinates}
    \lean{TNLean.PEPS.regularThreeCyclePlaquetteCoordinates}
    \leanok
    Let $J$ index group labels, and choose distinct $e_0,e_1,e_2\in J$.
    Define a permutation $C$ of $G^J$ which retains the other coordinates and
    sends the selected directed labels $(t_0,t_1,t_2)$ to
    \begin{align}
        C(t_0,t_1,t_2)
         & =(t_0^{-1}t_1,t_1^{-1}t_2,t_2^{-1}).
    \end{align}
    Its inverse sends $(a,b,c)$ to
    $(c^{-1}b^{-1}a^{-1},c^{-1}b^{-1},c^{-1})$.
\end{definition}

\begin{theorem}[Evaluation and equivariance of three plaquette coordinates]%
    \label{thm:peps_regular_three_plaquette_coordinates}
    \lean{TNLean.PEPS.regularThreeCyclePlaquetteCoordinates_apply,
        TNLean.PEPS.regularThreeCyclePlaquetteCoordinates_conjugation}
    \leanok
    \uses{def:peps_regular_three_plaquette_coordinates}
    The selected coordinates of $Cz$ are
    $(z_{e_0}^{-1}z_{e_1},z_{e_1}^{-1}z_{e_2},z_{e_2}^{-1})$, and
    $C(xzx^{-1})=xC(z)x^{-1}$ for every common conjugating element $x$.
\end{theorem}
\begin{proof}\leanok
    Evaluate the defining cases at the three distinct indices.
    Conjugation preserves multiplication and inversion, so it commutes with
    each selected coordinate formula and with every retained coordinate.
\end{proof}

\begin{definition}[Exchange of two group labels]%
    \label{def:peps_regular_two_cycle_exchange}
    \lean{TNLean.PEPS.regularTwoCycleExchange}
    \leanok
    \uses{def:peps_regular_two_cycle_conjugation}
    For distinct selected indices, let $H$ conjugate the second label by the
    first and then interchange the two coordinates. It sends
    $(a,b)$ to $(aba^{-1},a)$ and retains every other label.
\end{definition}

\begin{theorem}[Equivariance and the double exchange formula]%
    \label{thm:peps_regular_two_cycle_exchange}
    \lean{TNLean.PEPS.regularTwoCycleExchange_conjugation,
        TNLean.PEPS.regularTwoCycleExchange_apply,
        TNLean.PEPS.regularTwoCycleExchange_apply_other,
        TNLean.PEPS.regularTwoCycleExchange_twice_apply,
        TNLean.PEPS.regularTwoCycleExchange_twice_product,
        TNLean.PEPS.regularTwoCycleExchange_twice_classes}
    \leanok
    \uses{def:peps_regular_two_cycle_exchange,
        thm:peps_conjugacy_class_invariance,
        thm:peps_regular_two_cycle_conjugation}
    The permutation $H$ commutes with simultaneous conjugation.
    Its square retains the other coordinates and satisfies
    \begin{align}
        H^2(a,b)                 & =(abab^{-1}a^{-1},aba^{-1}), \\
        (H^2(a,b))_0(H^2(a,b))_1 & =ab.
    \end{align}
    The two output labels are respectively conjugate to $a$ by $ab$ and to
    $b$ by $a$.
\end{theorem}
\begin{proof}\leanok
    Evaluate the controlled conjugation followed by the interchange.
    A second evaluation gives the formula for $H^2$; multiplication cancels
    $b^{-1}a^{-1}ab$ and gives $ab$. Both component formulas commute with
    common conjugation.
\end{proof}

\begin{definition}[Oriented double exchange in three plaquette coordinates]%
    \label{def:peps_regular_three_plaquette_double_exchange}
    \lean{TNLean.PEPS.regularThreePlaquetteDoubleExchange,
        TNLean.PEPS.regularCycleOrientation,
        TNLean.PEPS.regularOrientedThreePlaquetteDoubleExchange}
    \leanok
    \uses{def:peps_regular_three_plaquette_coordinates,
        def:peps_regular_two_cycle_exchange}
    Define $\Phi=C^{-1}H^2C$ on directed bond labels.
    For a common orientation flag $\varepsilon\in\{0,1\}$, let
    $I_0z=z$ and $I_1z=z^{-1}$ coordinatewise, and define
    $\Phi_\varepsilon=I_\varepsilon^{-1}\Phi I_\varepsilon$.
\end{definition}

\begin{theorem}[Equivariance and action of oriented double exchange]%
    \label{thm:peps_regular_three_plaquette_double_exchange}
    \lean{TNLean.PEPS.regularThreePlaquetteDoubleExchange_conjugation,
        TNLean.PEPS.regularThreePlaquetteDoubleExchange_coordinates,
        TNLean.PEPS.regularCycleOrientation_conjugation,
        TNLean.PEPS.regularOrientedThreePlaquetteDoubleExchange_conjugation,
        TNLean.PEPS.regularOrientedThreePlaquetteDoubleExchange_coordinates,
        TNLean.PEPS.regularOrientedThreePlaquetteDoubleExchange_apply_other}
    \leanok
    \uses{thm:peps_regular_three_plaquette_coordinates,
        def:peps_regular_three_plaquette_double_exchange,
        thm:peps_regular_two_cycle_exchange}
    For every $z\in G^J$ and $x\in G$, the maps $I_\varepsilon$,
    $\Phi$ and $\Phi_\varepsilon$ commute with simultaneous conjugation.
    Their coordinate action satisfies
    \begin{align}
        C\Phi z                          & =H^2Cz,               \\
        CI_\varepsilon\Phi_\varepsilon z & =H^2CI_\varepsilon z.
    \end{align}
    Thus the selected plaquette coordinates change from $(a,b,c)$ to
    $((ab)a(ab)^{-1},aba^{-1},c)$, retaining all other coordinates.
    Their ordered product $abc$ and their individual conjugacy classes are
    preserved. At $(g,h,h^{-1})$, the last two output labels have joint product
    $ghg^{-1}h^{-1}$.
\end{theorem}
\begin{proof}\leanok
    Equivariance of $C$ implies equivariance of its inverse. Inversion is
    equivariant under conjugation, and so is $H$. Compose these identities
    in the defining order. The coordinate identities follow by cancelling
    the inverse coordinate maps; the double exchange formula gives the
    product and conjugacy assertions.
\end{proof}

\begin{theorem}[A fixed original-spin double exchange]%
    \label{thm:peps_regular_three_plaquette_double_exchange_physical}
    \lean{TNLean.PEPS.exists_unitary_regularOrientedThreePlaquetteDoubleExchange}
    \leanok
    \uses{thm:peps_regular_three_plaquette_double_exchange,
        thm:peps_regular_cycle_global_permutation}
    Let $G$ be finite, let $R$ be a finite graph region with a chosen spanning
    tree and root, and choose three distinct internal bonds outside the tree.
    Suppose the original site tensors are locally $G$-isometric for the
    regular incident-edge action. For each common orientation flag, there is
    one unitary $W$ on the original spins of $R$, chosen before every inserted
    cycle assignment $\omega$, boundary configuration $\theta$, and common
    exterior and crossing bond assignment $u$, such that
    \begin{align}
        W\Psi_R(u_\omega,\theta)
         & =\Psi_R(u_{\Phi_\varepsilon\omega},\theta), \\
        (W\otimes I_{R^c})\Psi(u_\omega,u)
         & =\Psi(u_{\Phi_\varepsilon\omega},u).
    \end{align}
    These are actual original-spin contractions. The operation is auxiliary to
    the conjugation and controlling-flux backreaction discussed in
    \cite[Section 6.6]{Schuch2010PEPS}. Identifying it with the prescribed
    oriented lattice string crossing remains separate.
\end{theorem}
\begin{proof}\leanok
    The permutation $\Phi_\varepsilon$ is conjugation-equivariant.
    Apply the actual physical and global cycle-permutation theorem to this
    single permutation. It supplies the same $W$ for every boundary column
    and its identity extension for every common exterior contraction.
\end{proof}
